% Complete single-arm results. No two-arm or unrestricted core claim.
\section{A sharp doubling bound for forward positive differences}
For an increasing nonempty arm write $F=\sum f_nz^n$ and put
$D_n=\max(f_n-f_{n-1},0)$ for $n\ge2$, with $D_0=D_1=0$.
\begin{theorem}
For every $n\ge3$,
\[
                            D_n\le2D_{n-1}.
\]
The constant two is sharp. Thus $D_n/2^n$ is nonincreasing
from rank two onward.
\end{theorem}
\begin{proof}
Normalize the first part to an initial run of ones, without
changing the polynomial. A rectangular arm has $D=0$.
Append a part $t\ge2$. If the old arm is rectangular, the
new positive differences are a terminal truncation of
$z^2[t-1]$, so the assertion holds.

Otherwise the old arm has $D_2=1$. Let $\kappa$ be its first
missing weight. Then $3\le\kappa\le t+1$: the old maximum
part is at most $t$, and the single-letter word $(t+1)$ is
missing. Put $G=(1-2z)D$. Induction says
$G_2=1$ and $G_i\le0$ for $i\ge3$. Moreover $G_\kappa\le-1$.
Indeed ranks below $\kappa$ are full, while at rank $\kappa$
at least one word is missing; therefore
$D_{\kappa-1}=2^{\kappa-3}$ and
$D_\kappa\le2^{\kappa-2}-1$.

Consider the nonnegative polynomial
\[
                         P=z^2[t-1]+z[t]D.
\]
Its coefficient of $(1-2z)P$ is nonpositive at every rank
$n\ge3$. For $3\le n\le t$, the first summand contributes
$-1$ and the window of $G$ contributes at most $1$.
At $n=t+1$ these bounds are $-2$ and $1$.
At $n=t+2$ the coefficient is
$\sum_{i=2}^{t+1}G_i\le0$, because this sum includes
$G_2=1$ and $G_\kappa\le-1$.
For $n>t+2$ the window includes only nonpositive $G_i$.

Let $m$ be the last mode of the old $F$. At ranks through
$m+1$, the new positive differences agree with $P$, since
$F_{\rm new}=1+z[t]F$ and all old differences through $m$
are nonnegative. Beyond $m+1$ the new positive differences
are nonincreasing: while positive they equal
$f_{n-1}-f_{n-t-1}$, with the first term nonincreasing and
the second nondecreasing. Once both indices exceed $m$,
this expression is nonpositive. This also checks the joining
edge at $m+1$. Consequently the doubling bound for $P$
gives the bound for the actual positive differences at every rank.

For sharpness, the increasing host $(1,2,3)$ contains every
word of weight at most three. Thus $D_2=1$ and $D_3=2$.
\end{proof}

The proof needs only that each appended part $t$ is at least
$\kappa-1$ for the preceding nonrectangular arm; monotonicity
of the host is a sufficient way to guarantee this condition.
The theorem concerns forward differences of one arm. The
analogous doubling bound for two-arm forward differences is
false, as recorded previously. The general weighted reverse
difference single-crossing conjecture therefore remains open.

The doubling bound is also a consequence of the earlier
single-arm deficit mode ordering. In fact, on the increasing
side of $F$, for $j=a-n+2$ one has
$\Delta R_j=D_n-2D_{n-1}$; this rank is on the decreasing
side of $R$. The preceding proof supplies an independent
universality-defect induction and is not claimed as a
logically independent new restriction beyond that ordering.

\begin{theorem}[Shape of the positive forward differences]
For every formal arm with an initial nonempty run of ones
and all subsequent parts at least two, in arbitrary order,
the sequence $(D_n)_{n\ge2}$ is unimodal.
\end{theorem}
\begin{proof}
Use induction on the number of appended parts. A rectangular
arm has zero positive differences. After the first part, the
positive differences form an interval of ones.

For the inductive step the old nonzero $D$ has initial support
$[2,e]$ and $D_2=1$. As above, form
$P=z^2[t-1]+W$, where $W=z[t]D$. The interval window $W$
is unimodal. At every rank $n\ge3$, its difference agrees
with that of $P$ except at $n=t+1$, where $P$ has one unit
subtracted. The new positive unit at rank two cannot create
a descent.

At the exceptional rank,
$\Delta W_{t+1}=D_t$. If $D_t\ge1$, the exceptional
difference is still nonnegative. If $D_t=0$, initial interval
support implies $D_j=0$ for all $j\ge t$, so
$\Delta W_n=D_{n-1}-D_{n-t-1}\le0$ for every $n\ge t+1$.
The exceptional rank is therefore on a nonincreasing tail.
In both cases $P$ is unimodal.

The actual new positive differences agree with $P$ through
one rank after the old last mode and are nonincreasing
afterward, by the crossing calculation in the preceding proof.
This splice preserves unimodality and completes induction.
\end{proof}

The initial support assertion used in this induction is also
preserved: $P$ is positive from rank two until its support
ends, and the terminal nonincreasing splice cannot insert
an internal zero before a later positive value.
